% [原文 p.83]
\chapter{算子}
\section{外幂算子}
\takeaway{本节把 K 理论中的算子定义为自然变换 $F_X\colon K(X)\to K(X)$，用外幂与对称幂构造 $\lambda_t$、$S_t$ 及其系数算子，再由变量替换得到 $\gamma_t$。核心结论是定理 3.1.7：$\gamma$ 算子生成的算子环同构于形式幂级数环。}

$K$ 理论中的算子 $F$，指的是自然变换 $F_X\colon K(X)\to K(X)$。也就是说，对每个空间 $X$ 都有一个（集合意义的）映射 $F_X\colon K(X)\to K(X)$，并且对任意连续映射 $f\colon X\to Y$，都有
\[
F_X f^*=f^* F_Y
\]
成立。

% [Tutor Note: 这里要求的是自然变换，而不是对各个空间任意指定的自同态：$F_X f^*=f^*F_Y$ 必须对一切连续映射 $f$ 成立。只在一个空间上给定的同态一般拼不成算子；本节构造的 $\\lambda_t$、$\\gamma_t$ 都自动满足这一自然性。]

设算子 $F$ 与 $G$ 满足：对任意可写成线丛之和的向量丛 $E$，以及任意整数 $n$，都有 $F([E]-n)=G([E]-n)$。此时，由上一章的分裂原理立刻可知，对所有 $x\in K(X)$ 都有 $F(x)=G(x)$。

利用外幂算子，可以用各种方式定义各种算子。先说明 Grothendieck 的方法。

设 $V$ 为空间 $X$ 上的向量丛，定义幂级数 $\lambda_t[V]\in K(X)[[t]]$ 为
\[
\sum_{i=0}^{\infty}t^i[\lambda^i(V)].
\]
由同构
\[
\lambda^k(V\oplus W)\cong\sum_{i+j=k}\lambda^i(V)\otimes\lambda^j(W),
\]
对任意向量丛 $V,W$ 得到
\[
\lambda_t[V\oplus W]=\lambda_t[V]\lambda_t[W].
\]
对任意 $W$，幂级数 $\lambda_t[W]$ 都是 $K(X)[[t]]$ 中的可逆元。因为它的常数项为 1。

这样，就得到从加法半群 $\Vect(X)$ 到 $K(X)$ 上常数项为 1 的幂级数所成乘法群的同态
% [原文 p.84]
\[
\lambda_t\colon \Vect(X)\longrightarrow 1+K(X)[[t]]^{+}.
\]
由 $K(X)$ 的普遍性，它唯一地扩张为同态
\[
\lambda_t\colon K(X)\longrightarrow 1+K(X)[[t]]^{+}.
\]
于是，取 $t^i$ 的系数，便得到算子
\[
\lambda^i\colon K(X)\longrightarrow K(X).
\]
具体地说，有
\[
\lambda_t([V]-[W])=\lambda_t[V]\lambda_t[W]^{-1}.
\]

用十分类似的方法，可以处理对称幂（symmetric power）$S^i(V)$。由于
\[
S^k(V\oplus W)\cong\sum_{i+j=k}S^i(V)\otimes S^j(W),
\]
便得到同态
\[
S_t\colon K(X)\longrightarrow 1+K(X)[[t]]^{+},
\]
其系数定义了算子
\[
S^i\colon K(X)\longrightarrow K(X).
\]
注意，当 $L$ 为线丛时，
\[
\lambda_t(L)=1+tL,\qquad
S_t(L)=1+tL+t^2L^2+\cdots=(1-tL)^{-1}.
\]
于是
% [原文 p.85]
\[
\lambda_{-t}(L)S_t(L)=1
\]
成立。因此，若 $V$ 是线丛之和，则 $\lambda_{-t}[V]S_t[V]=1$。所以对任意 $x\in K(X)$ 都有 $\lambda_{-t}(x)S_t(x)=1$，并且
\[
\lambda_t([V]-[W])=\lambda_t[V]S_{-t}[W]
\]
成立。也就是说，
\[
\lambda^k([V]-[W])=\sum_{i+j=k}(-1)^j\lambda^i[V]S^j[W].
\]
这给出了算子 $\lambda^i$ 作为向量丛上的算子的具体公式。

这里回想一下，对任意向量丛 $E$，$\dim E_x$ 是 $x$ 的局部常值函数。由于已假定 $X$ 是紧空间，
\[
\dim E=\sup_{x\in X}\dim E_x
\]
是有限的。外幂有一个基本性质：若 $i>\dim E$，则
\[
\lambda^i(E)=0
\]
成立。$K(X)$ 的元素，当它以真正的向量丛为代表元时，也就是当它落在 $\Vect(X)$ 的像中时，称为正的（positive），写作 $x\ge 0$。此时，
\[
x\ge 0\ \Longrightarrow\ \lambda_t(x)\in K(X)[t]
\]
成立。

% [Tutor Note: 此处的「正」只是说 $x$ 落在 $\\Vect(X)$ 的像中，并非来自 $K(X)$ 的任何序结构；$K(X)$ 上没有与 $\\oplus$ 相容的全序，切勿把 $x\\ge 0$ 读成实数大小。]

许多问题中重要的不是 $\dim E$，而是如下定义的另一整数。首先，取向量丛 $\rank E_x$，使其各点的纤维维数等于 $\dim E_x$，且在 $X$ 的每个连通分支上都是平凡丛。若 $X$ 连通，则 $\rank E$ 不外乎是维数等于 $\dim E$ 的平凡丛。此时，$E\mapsto\rank E$ 诱导（幂等）环自同态
% [原文 p.86]
\[
\rank\colon K(X)\longrightarrow K(X).
\]
这常被称为增广映射（augmentation）。这个自同态的核是理想，记作 $K_1(X)$。对于带基点的连通空间，显然
\[
K_1(X)=\tilde K(X).
\]
对任意 $x\in K(X)$，有
\[
x-\rank x\in K_1(X).
\]

% [Tutor Note: 原日译本 p.85 此处印作 $\\rank E_x$，同一段下文与全节各处则作 $\\rank E$。这里的 $\\rank E$ 是一个丛（各连通分支上纤维维数为 $\\dim E_x$ 的平凡丛），不是秩函数，下标 $x$ 疑为衍字；后文 $E<本地路径> E$ 才是那个环自同态。]

这里，对任意 $x\in K(X)$，定义 $\dim_K x$ 为满足条件
\[
x-\rank x+n\ge 0
\]
的最小整数 $n$。$K(X)$ 的任意元素都可以对某个向量丛 $V$ 写成 $[V]-n$ 的形式，所以可知 $\dim_K x$ 对一切 $x\in K(X)$ 都有限。对向量丛 $E$，显然
\[
\dim_K[E]\le\dim E
\]
成立。令 $x_1=x-\rank x$，则
\[
\dim_K x=\dim_K x_1,
\]
所以要注意，本质上 $\dim_K$ 是 $K(X)$ 模 $K_1(X)$ 的商上的函数。

这里引入新的算子 $\gamma^i$ 较为方便。$\gamma^i$ 与 $\dim_K$ 的关系，将和 $\lambda^i$ 与向量丛维数的关系相同。仍按 Grothendieck 的做法，定义
\[
\gamma_t(x)=\lambda_{t/(1-t)}(x)\in K(X)[[t]],
\]
则 $\gamma_t(x+y)=\gamma_t(x)\gamma_t(y)$。因此，对每个 $i$，得到算子
% [原文 p.87]
\[
\gamma^i\colon K(X)\longrightarrow K(X).
\]
若令 $s=t/(1-t)$、$t=s/(1+s)$，则
\[
\lambda_s(x)=\gamma_{s/(1+s)}(x)
\]
成立，所以 $\gamma^i$ 是 $\lambda^j$（$j\le i$）的线性组合，反之亦然。

% [Tutor Note: $\\gamma_t$ 只是 $\\lambda_t$ 的变量替换，并不产生新的信息；改用 $\\gamma^i$ 的理由在命题 3.1.1：$\\gamma$-类在 $K_1(X)$ 上幂零，因而形式幂级数可作用于任意 $x\\in K_1$，这正是定理 3.1.7 的出发点。]

要注意，
\[
\gamma_t(1)=(1-t)^{-1},
\]
并且对线丛 $L$，有
\[
\gamma_t([L]-1)=1+t([L]-1).
\]

\begin{proposition*}{3.1.1}
设 $x\in K_1(X)$。则 $\gamma_t(x)$ 是次数不超过 $\dim_K x$ 的多项式。
\end{proposition*}

\begin{proof}
令 $n=\dim_K x$，则 $x+n\ge 0$。因此，对某个向量丛 $E$ 有 $x+n=[E]$。又 $\dim E=n$，故对 $i>n$ 有
\[
\lambda^i(E)=0
\]
成立。因此 $\lambda_t(x+n)$ 是次数不超过 $n$ 的多项式。现在，
\begin{align*}
\gamma_t(x)&=\gamma_t(x+n)\gamma_t(1)^{-n}\\
&=\lambda_{t/(1-t)}(x+n)(1-t)^n\\
&=\sum_{i=0}^{n}\lambda^i(x+n)t^i(1-t)^{n-i},
\end{align*}
所以如命题所述，$\gamma_t(x)$ 是次数不超过 $n$ 的多项式。
\end{proof}

% [Tutor Note: 证明中「$\\dim E=n$」省了一步：由 $x+n=[E]$ 与 $\\rank x=0$ 得 $\\rank E=n$，即 $E$ 在各连通分支上的纤维维数是 $n$，故 $\\dim E=n$；否则后面 $\\lambda^i(E)=0$（$i>n$）不能直接用。]

% [原文 p.88]
这里，定义 $\dim_\gamma x$ 为满足 $\gamma^n(x-\rank x)\ne 0$ 的最大整数 $n$，并令
\begin{gather*}
\dim_K X=\sup_{x\in K(X)}\dim_K x,\\
\dim_\gamma X=\sup_{x\in K(X)}\dim_\gamma x.
\end{gather*}
由命题 3.1.1，
\[
\dim_\gamma x\le\dim_K x,\qquad
\dim_\gamma X\le\dim_K X
\]
成立。

下面在较宽松的限制下证明 $\dim_K X$ 有限。为此，需要准备几个关于对称函数的引理。

\begin{lemma*}{3.1.2}
设 $x_1,\dots,x_n$ 为未定元。此时，$\Z[x_1,\dots,x_n]$ 中次数大于 $n(n-1)$ 的任意齐次多项式，都属于由 $(x_1,\dots,x_n)$ 的正次数对称函数生成的理想。
\end{lemma*}

\begin{proof}
设 $\sigma_i(x_1,\dots,x_n)$ 为 $i$ 次初等对称多项式。此时，方程
\[
x^n-\sigma_1x^{n-1}+\sigma_2x^{n-2}+\cdots+(-1)^n\sigma_n=0
\]
以 $x=x_i$ 为根。因此，$x_i^n$ 属于由 $\sigma_1,\dots,\sigma_n$ 生成的理想。然而，$x_1,\dots,x_n$ 的次数大于 $n(n-1)$ 的单项式被某个 $i$ 对应的 $x_i^n$ 整除，故属于这个理想。
\end{proof}

% [Tutor Note: 界 $n(n-1)$ 来自「$n$ 个变量、每个指数不超过 $n-1$」时可能达到的最高次数；证明中「被某个 $x_i^n$ 整除」用的是抽屉原理：次数一旦超过 $n(n-1)$，必有某个指数不小于 $n$。]

\begin{lemma*}{3.1.3}
设 $x_1,\dots,x_n,y_1,\dots,y_m$ 为未定元，
\[
a_i=\sigma_i(x_1,\dots,x_n),\qquad
b_i=\sigma_i(y_1,\dots,y_m)
\]
为初等对称多项式。设 $I$ 为 $\Z[a,b]$ 的任意理想，$J$ 为它在 $\Z[x,y]$ 中的扩张。此时
\[
J\cap\Z[a,b]=I
\]
成立。
\end{lemma*}

% [原文 p.89]
\begin{proof}
众所周知，$\Z[x]$ 是以单项式
\[
x^r=x_1^{r_1}x_2^{r_2}\cdots x_{n-1}^{r_{n-1}}\qquad(0\le r_i\le n-i)
\]
为基的自由 $\Z[a]$ 模。因此，$\Z[x,y]=\Z[x]\otimes\Z[y]$ 是 $\Z[a,b]=\Z[a]\otimes\Z[b]$ 上的自由模，以单项式 $x^ry^s$ 为基。此时，理想 $J\subset\Z[x,y]$ 由一切形如
\[
f=\sum f_{r,s}x^ry^s\qquad(f_{r,s}\in I)
\]
的元素 $f$ 组成。由于 $x^ry^s$ 是自由基，$f$ 属于 $\Z[a,b]$ 当且仅当对 $(r,s)\ne(0,0)$ 有 $f_{r,s}=0$。此时
\[
f=f_{0,0}\in I.
\]
因此 $J\cap\Z[a,b]=I$。
\end{proof}

\begin{remark}{}
这个引理断言的是 $\Z[a,b]/I$ 到 $\Z[x,y]/J$ 的嵌入，本质上是代数的分裂原理。不过，这个引理只是形式性质。处理代数的、形式的结果时，与其用位相的分裂原理（推论 2.7.11），不如用这个引理更好。位相的分裂原理当然基于周期性定理这一深刻结果，只应在处理同样深刻的性质时使用。
\end{remark}

% [Tutor Note: 作者在此明确划出两条路线：引理 3.1.3 是「代数的分裂原理」，只依赖对称函数恒等式；§2.7 的位相分裂原理依赖 Bott 周期性。若用后者去证本节的形式结论，就会绕回第 2 章的深刻定理，形成循环论证。]

\begin{lemma*}{3.1.4}
设 $K$ 为（以 1 为单位元的）交换环。设 $K[t]$ 的元素
\begin{align*}
a(t)&=1+a_1t+a_2t^2+\cdots+a_nt^n,\\
b(t)&=1+b_1t+b_2t^2+\cdots+b_mt^m
\end{align*}
满足 $a(t)b(t)=1$。此时存在整数 $N=N(m,n)$，使得任意单项式
\[
a_1^{r_1}a_2^{r_2}\cdots a_n^{r_n}
\]
只要其重量 $\sum_j jr_j>N$，则为零。
\end{lemma*}

% [原文 p.90]
\begin{proof}
考虑普遍情形即可：设 $a_1,\dots,a_n,b_1,\dots,b_m$ 为未定元。此时只需证明：$a_i$ 的单项式 $\alpha$ 只要重量不小于 $N$，就属于由
\[
c_k=\sum_{i+j=k}a_ib_j\qquad(k=1,\dots,mn,\ a_0=b_0=1)
\]
生成的理想 $I$。依引理 3.1.3，只需引入未定元 $x_1,\dots,x_n,y_1,\dots,y_m$，证明 $\alpha$ 属于扩张理想 $J$ 即可。然而，$c_k$ 恰好是 $(m+n)$ 个变量 $x_1,\dots,x_n,y_1,\dots,y_m$ 的 $k$ 次初等对称多项式。于是取 $N=(m+n)(m+n-1)$，应用引理 3.1.2，结论便成立。
\end{proof}

\begin{remark}{}
引理 3.1.4 的证明中得到的 $N(m,n)$ 不是最佳值。更详细的讨论可以证明，最佳值是 $mn$。
\end{remark}

那么，来应用这些代数结果。

\begin{proposition*}{3.1.5}
设 $x\in K_1(X)$。此时存在依赖于 $x$ 的整数 $N$，使得对任意单项式
\[
\gamma^{i_1}(x)\gamma^{i_2}(x)\cdots\gamma^{i_k}(x),
\]
只要其重量 $\sum_{j=1}^{k}i_j$ 大于 $N$，就等于零。
\end{proposition*}

\begin{proof}
对多项式 $\gamma_t(x)$、$\gamma_t(-x)$ 应用引理 3.1.4 即可。因此，$N$ 依赖于 $\dim_\gamma x$、$\dim_\gamma(-x)$。
\end{proof}

由 $\gamma^1(x)=x$，可导出下面的推论。

\begin{corollary*}{3.1.6}
任意 $x\in K_1(X)$ 都是幂零的。
\end{corollary*}

分别定义各 $\gamma^i$ 的次数（degree）为 1。此时，由 $\gamma^i$ 组成的任意单项式，其重量与次数满足
% [原文 p.91]
\[
\text{重量}\ \ge\ \text{次数}.
\]
因此，结合命题 3.1.5，若 $x\in K_1(X)$，则 $\gamma^i(x)$ 的单项式只要次数足够高就为零。于是，可以把 $\gamma^i$ 的形式幂级数（formal power series）作用于任意 $x\in K_1(X)$。以下用 $\operatorname{Op}(K_1,K)$ 表示从 $K_1$ 到 $K$ 的算子的全体。$\operatorname{Op}(K_1,K)$ 带有由 $K$ 的环结构诱导的环结构（值的加法与乘法）。此时，由以上所述，得到环同态
\[
\varphi\colon \Z[[\gamma^1,\dots,\gamma^n,\dots]]\longrightarrow \operatorname{Op}(K_1,K).
\]

\begin{theorem*}{3.1.7}
\[
\varphi\colon \Z[[\gamma^1,\dots,\gamma^n,\dots]]\longrightarrow \operatorname{Op}(K_1,K)
\]
是同构映射。
\end{theorem*}

\begin{proof}
设 $Y_{n,m}$ 为 $n$ 个 $P_m(\C)$ 的乘积。利用 $P_m(\C)$ 的基点 $P_0(\C)$，$Y_{n,m}$ 连同对 $n'\ge n$、$m'\ge m$ 的包含映射
\[
Y_{n,m}\longrightarrow Y_{n',m'}
\]
构成空间的归纳系。此时，$K(Y_{n,m})$ 是群的射影系，并且
\begin{gather*}
K(Y_{n,m})=\Z[x_1,\dots,x_n]/(x_1^{m+1},\dots,x_n^{m+n}),\\
\lim_{m}K(Y_{n,m})=\Z[[x_1,\dots,x_n]],\\
\lim_{m,n}K(Y_{n,m})=\lim_{n}\Z[[x_1,\dots,x_n]]
\end{gather*}
成立。任意算子都诱导射影极限上的算子。因此，可以定义映射
% [原文 p.92]
\[
\psi\colon \operatorname{Op}(K_1,K)\longrightarrow \lim_{n}\Z[[x_1,\dots,x_n]]
\]
为
\[
\psi(f)=\lim_{n}f(x_1+x_2+\cdots+x_n).
\]
在 $K(Y_{n,m})$ 中
\[
\gamma_t(x_1+x_2+\cdots+x_n)=\prod_{i=1}^{n}(1+x_it)
\]
成立，所以若 $\sigma_i$ 为 $i$ 次初等对称多项式，则
\[
\psi\circ\varphi(\gamma^i)=\lim_{n}\sigma_i(x_1,\dots,x_n)
\]
成立。因此，特别地 $\psi\circ\varphi$ 是单射，从而 $\varphi$ 也是单射。进一步，$\psi\circ\varphi$ 的像是
\[
\Z[[\sigma_1,\dots,\sigma_n]],
\]
它与
\[
\lim_{n}\Z[[x_1,\dots,x_n]]^{S_n}
\]
相同。这里 $[\cdot]^{S_n}$ 表示在对称群 $S_n$ 下不变的部分环。然而，对一切 $f\in\operatorname{Op}(K_1,K)$，
\[
\psi(f)=\lim_{n}f(x_1+\cdots+x_n)
\]
都属于这个环。换言之，
\[
\Image(\psi\circ\varphi)=\Image\psi.
\]
为完成这一证明，只需说明 $\psi$ 是单射。于是假设 $\psi(f)=0$。空间 $X$ 上的任意线丛都同构于由到 $P_n(\C)$ 的映射拉回的万有丛，所以当 $E$ 是 $n$ 个线丛之和时，总有
\[
f([E]-n)=0.
\]
由这一点，依分裂原理，对所有 $x\in K_1$ 有 $f(x)=0$ 成立。也就是说，$f$ 正是所求的零算子。
\end{proof}

% [Tutor Note: $K(Y_{n,m})$ 一行中的 $x_i$ 不是超平面线丛的类本身，而是 $[H_i]-1$；正因如此才有 $\\gamma_t(x_1+\\cdots+x_n)=\\prod(1+x_it)$。若误把 $x_i$ 当作 $[H_i]$，这条公式会差一个因子。]

% 块尾跨页句：このことから、分裂原理によって、すべての $x\in K_1$ に対して …（已写完至 p.93「$f(x)=0$ が導かれる／すなわち、f は求めるとおりゼロ作用素である」；p.93 其余正文由 ch03b 接续，页标记「原文 p.93」由 ch03b 写）
% CHUNK ch03a DONE
